% !TEX program = pdflatex
% Theory document for the Krylov-subspace Model-Order Reduction (MOR) fast
% frequency sweep module of the BP-FEM solver. Self-contained; compile with
% `pdflatex theory.tex`.

\documentclass[11pt,a4paper]{article}

\usepackage[utf8]{inputenc}
\usepackage[T1]{fontenc}
\usepackage{amsmath,amssymb,amsthm,bm}
\usepackage{mathtools}
\usepackage{booktabs}
\usepackage{geometry}
\usepackage{hyperref}
\usepackage{algorithm}
\usepackage{algpseudocode}
\usepackage{xcolor}

\geometry{margin=2.5cm}
\hypersetup{colorlinks=true,linkcolor=blue!50!black,urlcolor=blue!50!black,citecolor=blue!50!black}

\newcommand{\R}{\mathbb{R}}
\newcommand{\C}{\mathbb{C}}
\newcommand{\im}{\mathrm{i}}
\DeclareMathOperator{\range}{range}
\DeclareMathOperator{\diag}{diag}
\DeclareMathOperator{\Real}{Re}
\DeclareMathOperator{\Imag}{Im}
\DeclareMathOperator{\spanop}{span}

\newtheorem{definition}{Definition}
\newtheorem{theorem}{Theorem}
\newtheorem{remark}{Remark}

\title{Krylov-Subspace Model-Order Reduction \\ for Fast Frequency Sweeps in the BP-FEM Solver}
\author{BP-FEM Optimization Track}
\date{Draft for review}

\begin{document}
\maketitle
\tableofcontents
\bigskip

%---------------------------------------------------------------
\section{Scope}

This document specifies the mathematical foundation of the Krylov-subspace
Model-Order Reduction (MOR) module to be added to the BP-FEM solver as a
\emph{fast frequency sweep} option, alongside the existing point-by-point
direct solver and the BiCGStab fallback. The goal is to reduce a sweep of
$N_f$ frequency points over a band of interest from $\mathcal{O}(N_f)$
factorizations of the full system to a single (or few) factorization(s)
plus $\mathcal{O}(N_f)$ evaluations on a reduced model whose dimension is
two orders of magnitude smaller.

The document is intentionally implementation-aware: every symbol below has
a matching counterpart in the existing C++ source tree
(\texttt{FEMAssembler}, \texttt{PortModeSolver},
\texttt{MklPardisoSolver}, \texttt{ResultExtractor}, \texttt{SparseMatrix}).
Where convenient, source-level names are noted in
\textcolor{gray}{gray} comments such as
\textcolor{gray}{\texttt{// FEMAssembler::assemble}}.

%---------------------------------------------------------------
\section{Continuous problem and FE discretization}

\subsection{Strong form}

We solve the time-harmonic curl--curl problem in a bounded Lipschitz
domain $\Omega \subset \R^3$ partitioned into wave-port faces
$\{\Gamma_p\}_{p=1}^{N_p}$ and a perfect electric conductor (PEC) wall
$\Gamma_{\mathrm{PEC}}$ with $\partial\Omega = \Gamma_{\mathrm{PEC}}\cup
\bigcup_p \Gamma_p$. The time convention is $\mathrm{e}^{+\im\omega t}$.
For an angular frequency $\omega = 2\pi f$ the electric field
$\bm E(\bm r,\omega) \in H(\mathrm{curl};\Omega)$ satisfies
\begin{align}
  \nabla\times\!\big(\mu_r^{-1}\nabla\times\bm E\big)
  -k_0^2\,\varepsilon_c(\omega)\,\bm E
  &= \bm 0
  & &\text{in } \Omega,\label{eq:strong}\\
  \hat{\bm n}\times\bm E &= \bm 0
  & &\text{on } \Gamma_{\mathrm{PEC}},\\
  \hat{\bm n}\times\!\big(\mu_r^{-1}\nabla\times\bm E\big)
  &= \mathcal B_p(\bm E;\omega) + \bm g_p(\omega)
  & &\text{on } \Gamma_p,
\end{align}
where $k_0 = \omega/c_0$, $c_0$ is the vacuum speed of light, and
$\varepsilon_c(\omega) = \varepsilon_r - \im\sigma/(\omega\varepsilon_0)$
encodes optional volumetric losses
\textcolor{gray}{(\texttt{FEMAssembler::assemble})}.

\subsection{Wave-port boundary form}

Each port $\Gamma_p$ admits a numerically computed dominant transverse
mode $\bm e_p^{(t)}: \Gamma_p \to \R^3$ with cutoff wavenumber
$k_{c,p}$ and propagation constant
\(
  \beta_p(\omega) = \sqrt{k_0^2 - k_{c,p}^2}
\),
absorbing the propagating mode and exciting the projector
\(
  \bm g_p = -2\,\im\beta_p\,a_p^{\mathrm{inc}}\,(\hat{\bm n}\times\bm e_p^{(t)})
\)
where $a_p^{\mathrm{inc}}$ is the complex amplitude of the incident wave
in $\sqrt{\mathrm{W}}$ units. The L\textsuperscript{2}-shape of
$\bm e_p^{(t)}$ is normalized to
$\int_{\Gamma_p} |\bm e_p^{(t)}|^2\,\mathrm dS = 1$
and the power normalization $s_p(\omega)$ is supplied by
\textcolor{gray}{\texttt{powerNormalizationFactor}}:
\begin{equation}
  s_p(\omega) = \frac{1}{\sqrt{P_p^{\mathrm{unit}}(\omega)}}, \qquad
  P_p^{\mathrm{unit}}(\omega) = \frac{\beta_p(\omega)}{2\omega\mu_0}
  \int_{\Gamma_p}|\bm e_p^{(t)}|^2\,\mathrm dS.
  \label{eq:powernorm}
\end{equation}
With this normalization the incident amplitude carried by the right-hand
side becomes $\sqrt{P_p^{\mathrm W}}\,s_p(\omega)$ where
$P_p^{\mathrm W}$ is the requested incident power
\textcolor{gray}{(\texttt{PortDefinition::magnitudeW})}.

\subsection{Galerkin discretization}

Let $V_h\subset H(\mathrm{curl};\Omega)$ be the finite-element subspace
spanned by H(curl)-conforming Nedelec/Whitney edge elements
(0\textsuperscript{th} or 1\textsuperscript{st}-order hierarchical;
\textcolor{gray}{\texttt{EdgeTopology}}). Denote the global degree-of-freedom
vector $\bm x(\omega)\in\C^{n}$. The discrete bilinear forms and load
vector at frequency $\omega$ read
\begin{align}
  \bm K_{ij} &= \int_\Omega \mu_r^{-1}\,(\nabla\times\bm w_i)\cdot
                 (\nabla\times\bm w_j)\,\mathrm dV,
  & \text{[stiffness]}\\
  \bm M_{ij} &= \int_\Omega \varepsilon_c(\omega)\,\bm w_i\cdot\bm w_j\,\mathrm dV,
  & \text{[mass; weakly $\omega$-dependent if $\sigma\!\neq\!0$]}\\
  \bm m_p    &= \int_{\Gamma_p}(\hat{\bm n}\times\bm e_p^{(t)})\cdot\bm w_i\,\mathrm dS,
  & \text{[port projector / coupling vector]}\\
  \bm f(\omega) &= \sum_{p=1}^{N_p} 2\,\im\beta_p(\omega)\,\sqrt{P_p^{\mathrm W}}\,s_p(\omega)\,\delta_p^{\mathrm{exc}}\,\bm m_p,
  & \text{[wave-port RHS]}
\end{align}
which assemble into the discrete system
\begin{equation}
  \boxed{\;
  \bm A(\omega)\,\bm x(\omega) = \bm f(\omega),\qquad
  \bm A(\omega) = \bm K - k_0^2\,\bm M
  + \im\sum_{p=1}^{N_p} \beta_p(\omega)\,\bm m_p \bm m_p^{\!\top}.
  \;}
  \label{eq:discrete}
\end{equation}
Equation~\eqref{eq:discrete} is the \emph{exact} system the existing
solver assembles per frequency
\textcolor{gray}{(\texttt{FEMAssembler::assemble} for $\bm K,\bm M,\bm m_p$;
\texttt{applyWavePorts} for the rank-one port term).}
The S-parameter postprocessor recovers
\begin{equation}
  b_p(\omega) = \frac{\bm m_p^{\!\top}\,\bm x(\omega)}{s_p(\omega)},
  \qquad
  S_{pq}(\omega) = \frac{b_p(\omega)-a_p^{\mathrm{inc}}\,\delta_{pq}}{a_q^{\mathrm{inc}}}
  \quad\text{(for unit excitation $a_q^{\mathrm{inc}}=1$).}
  \label{eq:sparam}
\end{equation}

\begin{remark}[Why this form is MOR-friendly]
The system $\bm A(\omega)$ in~\eqref{eq:discrete} is a low-order
\emph{rational matrix function} of $\omega$:
\begin{itemize}
  \item $\bm K - k_0^2\bm M$ is a polynomial of degree~2 in $k_0$, hence in
        $\omega$, modulo a possibly weak $\omega$-dependence of
        $\varepsilon_c$ that is exactly handled by the model
        in~\S\ref{sec:dispersion}.
  \item Each port term $\im\beta_p(\omega)\bm m_p\bm m_p^{\!\top}$ is a
        rank-one update with frequency dependence concentrated in the scalar
        $\beta_p(\omega)=\sqrt{k_0^2-k_{c,p}^2}$.
\end{itemize}
This structure is what allows a Krylov subspace built once at one (or a
few) expansion points to capture the entire transfer function~$S(\omega)$.
\end{remark}

%---------------------------------------------------------------
\section{Linear time-invariant transfer model}

\subsection{State-space embedding}

To keep the algebra clean we treat $k_0^2$ (or, for second-order moment
matching, $\omega$) as the expansion parameter and write the parametric
system in descriptor form
\begin{equation}
  \big[\bm K - k_0^2\,\bm M + \im\sum_p \beta_p(\omega)\,\bm m_p\bm m_p^{\!\top}\big]\,
  \bm x = \bm B\,\bm u,
  \qquad \bm y = \bm C\,\bm x,
  \label{eq:lti}
\end{equation}
with input matrix $\bm B = [2\im\beta_1 s_1 \bm m_1\;\cdots\;2\im\beta_{N_p} s_{N_p}\bm m_{N_p}]$,
input vector $\bm u\in\C^{N_p}$ holding the per-port incident amplitudes,
and output matrix $\bm C = [\bm m_1\;\cdots\;\bm m_{N_p}]^{\!\top}/\diag(s_p)$.
The transfer function of interest is
\begin{equation}
  \bm H(\omega) = \bm C\,\bm A(\omega)^{-1}\bm B \in \C^{N_p\times N_p},
\end{equation}
whose $(p,q)$ entry is precisely $b_p/a_q^{\mathrm{inc}}$ from
\eqref{eq:sparam} up to a known additive identity for the reflection
diagonal.

\subsection{Two MOR formulations}

We support two formulations because they offer different tradeoffs and
the choice is not premature optimization:

\paragraph{Form A --- single-point ``shifted-and-inverted'' Krylov.}
Pick an expansion point $\omega_0$ (band-center). Define the shifted
operator
\(
  \bm A_0 = \bm A(\omega_0).
\)
The Krylov subspace
\begin{equation}
  \mathcal{K}_q(\bm A_0^{-1}\bm M,\,\bm A_0^{-1}\bm B)
  = \spanop\!\big\{\bm A_0^{-1}\bm B,\ (\bm A_0^{-1}\bm M)\bm A_0^{-1}\bm B,\ \ldots,
                    (\bm A_0^{-1}\bm M)^{q-1}\bm A_0^{-1}\bm B\big\}
  \label{eq:krylov-singlepoint}
\end{equation}
matches $2q$ moments of $\bm H$ around $\omega_0$ (block Arnoldi version,
PRIMA-style).

\paragraph{Form B --- multi-point (rational) Krylov.}
Choose a small set of expansion points
$\{\omega_0^{(1)},\dots,\omega_0^{(L)}\}$ inside the sweep band. For each
expansion point factorize $\bm A_0^{(\ell)}$ once and append
$q_\ell$ columns of the corresponding shift-and-invert Krylov subspace.
The union spans a rational Krylov subspace that matches moments at every
expansion point and is far more accurate over wide bands than Form~A.

Both forms project $\bm A(\omega)$ onto the subspace $\bm V$ to produce
the reduced model
\begin{equation}
  \widetilde{\bm A}(\omega) = \bm V^{\!*}\bm A(\omega)\bm V,\qquad
  \widetilde{\bm B} = \bm V^{\!*}\bm B,\qquad
  \widetilde{\bm C} = \bm C\,\bm V,
\end{equation}
yielding the order-$q$ reduced transfer function
$\widetilde{\bm H}(\omega) = \widetilde{\bm C}\,\widetilde{\bm A}(\omega)^{-1}\widetilde{\bm B}$.

\begin{theorem}[Moment-matching, single-point form]
Let $\bm V\in\C^{n\times q}$ be an orthonormal basis of
$\mathcal{K}_q(\bm A_0^{-1}\bm M,\bm A_0^{-1}\bm B)$. The Galerkin-type
reduced model with $\widetilde{\bm A}(\omega) = \bm V^{\!*}\bm A(\omega)\bm V$
matches the first $2q$ moments of the original transfer function
$\bm H(\omega)$ around $\omega_0$ (Grimme's theorem). In particular,
\(
  \tilde H_{pq}^{(k)}(\omega_0) = H_{pq}^{(k)}(\omega_0)
\)
for all $k=0,1,\dots,2q-1$ and $p,q$.
\end{theorem}

\begin{remark}[Symmetry preservation]
$\bm K,\bm M$ are real symmetric and $\bm m_p\bm m_p^{\!\top}$ are SPSD
in this assembly. With Galerkin projection ($\bm V^{\!*}$ on the left,
$\bm V$ on the right) and $\bm V$ orthonormal, the reduced system retains
the same complex-symmetric structure, so the reduced $\widetilde{\bm A}$
is a tiny dense complex-symmetric matrix solvable by LDL\textsuperscript{T}
in $\mathcal{O}(q^3)$.
\end{remark}

\subsection{Frequency-dependent material model}\label{sec:dispersion}

Let $\varepsilon_c(\omega)$ be a rational function of $\omega$. The
generic bilinear-form parametric expansion is
\begin{equation}
  \bm A(\omega) = \bm K + \sum_{i=0}^{N_a} \alpha_i(\omega)\,\bm A_i,\qquad
  \alpha_i(\omega)\in\{\,1,\;k_0^2,\;\im k_0,\;\im\beta_p(\omega),\ldots\,\}.
\end{equation}
Each $\bm A_i$ is assembled \emph{once} during the offline phase and
stored in CSR alongside the existing
\textcolor{gray}{\texttt{SparsePattern}} cache. The reduced operator becomes
\begin{equation}
  \widetilde{\bm A}(\omega) = \bm V^{\!*}\bm K\bm V + \sum_i \alpha_i(\omega)\,
                              \bm V^{\!*}\bm A_i\bm V,
\end{equation}
so the online cost per frequency is $\mathcal{O}(N_a q^2)$ assembly plus
$\mathcal{O}(q^3)$ solve, independent of the full DOF count $n$.

%---------------------------------------------------------------
\section{Algorithms}

\subsection{Block-Arnoldi for single-point form}

\begin{algorithm}[h]
\caption{Block-Arnoldi (single-point Krylov MOR)}
\label{alg:block-arnoldi}
\begin{algorithmic}[1]
\Require expansion point $\omega_0$, max moments $q$, residual tolerance $\tau$
\Statex Operators: $\bm K$, $\bm M$, $\bm B$, $\bm C$ already assembled
\State Form $\bm A_0 \gets \bm A(\omega_0)$ and factorize once via PARDISO
\Comment{\textcolor{gray}{\texttt{MklPardisoSolver}: re-use symbolic factorization across frequencies for the reduction phase}}
\State $\bm V_1 \gets \mathrm{orth}(\bm A_0^{-1} \bm B)$
\For{$j = 1,\dots,q$}
  \State $\bm W \gets \bm A_0^{-1}\,\bm M\,\bm V_j$
  \For{$i = 1,\dots,j$}
    \State $\bm H_{ij}\gets \bm V_i^{\!*}\bm W$,\quad $\bm W\gets \bm W - \bm V_i\,\bm H_{ij}$ \Comment{Modified Gram-Schmidt}
  \EndFor
  \State (Re-orthogonalize $\bm W$ against $\bm V_{1:j}$ if needed.)
  \State $\bm V_{j+1} \gets \mathrm{orth}(\bm W)$;\quad break if $\|\bm W\| < \tau\|\bm V_1\|$
\EndFor
\State $\bm V \gets [\bm V_1,\bm V_2,\ldots]$
\State Build reduced matrices: $\widetilde{\bm K}\!=\!\bm V^{\!*}\bm K\bm V$, $\widetilde{\bm M}\!=\!\bm V^{\!*}\bm M\bm V$, $\widetilde{\bm m}_p\!=\!\bm V^{\!*}\bm m_p$
\State \Return $\bm V$, $\widetilde{\bm K}$, $\widetilde{\bm M}$, $\{\widetilde{\bm m}_p\}$
\end{algorithmic}
\end{algorithm}

\subsection{Multi-point (rational) Krylov}

\begin{algorithm}[h]
\caption{Adaptive rational Krylov MOR sweep}
\label{alg:rational-krylov}
\begin{algorithmic}[1]
\Require sweep band $[\omega_a,\omega_b]$, target tolerance $\varepsilon_S$, max ROM order $q_{\max}$
\State Pick initial expansion points $\Sigma_0$ (e.g.\ $\{\omega_a,\omega_m,\omega_b\}$)
\State $\bm V \gets [\,]$
\For{each $\omega_0^{(\ell)} \in \Sigma_0$}
  \State Factorize $\bm A_0^{(\ell)} = \bm A(\omega_0^{(\ell)})$
  \State Run Algorithm~\ref{alg:block-arnoldi} for $q_\ell$ steps; append columns to $\bm V$
\EndFor
\State Orthonormalize $\bm V$ globally; build reduced operators
\Repeat
  \State Evaluate $\widetilde{\bm H}(\omega)$ on a probe set $\Omega_p\subset[\omega_a,\omega_b]$
  \State Compute residual indicator $\rho(\omega)$ from \S\ref{sec:residual}
  \State $\omega^\star \gets \arg\max_{\omega\in\Omega_p}\rho(\omega)$
  \If{$\rho(\omega^\star) > \varepsilon_S$ and $\dim\bm V < q_{\max}$}
    \State Factorize $\bm A(\omega^\star)$, append Krylov block, re-orthonormalize
  \EndIf
\Until{$\rho(\omega^\star) \le \varepsilon_S$ or budget exhausted}
\State \Return $\bm V$ and reduced operators
\end{algorithmic}
\end{algorithm}

%---------------------------------------------------------------
\section{Error and stability indicators}\label{sec:residual}

We expose three indicators; the user-facing tolerance is on (R3) but
(R1) and (R2) are essential for diagnostics.

\paragraph{(R1) Algebraic residual.} For an evaluated frequency the relative
residual of the reduced solution lifted back to the full space is
\begin{equation}
  \rho_{\mathrm{alg}}(\omega) =
  \frac{\|\bm A(\omega)\,\bm V\,\widetilde{\bm x}(\omega) - \bm f(\omega)\|_2}
       {\|\bm f(\omega)\|_2}.
\end{equation}
If~$\rho_{\mathrm{alg}}$ explodes the basis is too small or the expansion
point is poorly placed. This indicator requires only one SpMV with the
full system, no factorization.

\paragraph{(R2) Heuristic moment residual.} Fully embedded in the ROM:
\begin{equation}
  \rho_{\mathrm{mom}}(\omega) =
  \|\widetilde{\bm A}(\omega)\,\widetilde{\bm x}(\omega) - \widetilde{\bm B}\,\bm u\|_2
  / \|\widetilde{\bm B}\,\bm u\|_2.
\end{equation}
This is essentially a self-consistency check on the reduced linear solve;
useful for catching condition-number blowup of $\widetilde{\bm A}$.

\paragraph{(R3) S-parameter discrepancy at validation points.}
Pick a small validation set $\Omega_v\subset[\omega_a,\omega_b]$ and run the
existing direct solver. Define
\begin{equation}
  \rho_S = \max_{\omega\in\Omega_v} \big\|\widetilde{\bm S}(\omega)-\bm S(\omega)\big\|_F,
\end{equation}
which is the user contract: the ROM output must agree with the exact
solver to within $\rho_S \le \varepsilon_S$ (default 0.01 dB on every
$|S|$).

%---------------------------------------------------------------
\section{Output reconstruction}

\subsection{S-parameter evaluation}

For each requested $\omega$:
\begin{enumerate}
  \item Solve the dense reduced system
        $\widetilde{\bm A}(\omega)\,\widetilde{\bm x}(\omega) = \widetilde{\bm B}\,\bm u$
        with the (cached) LDL\textsuperscript{T} of $\widetilde{\bm A}(\omega)$.
  \item Form $\bm y(\omega) = \widetilde{\bm C}\,\widetilde{\bm x}(\omega)$.
  \item Apply the existing power normalization
        $b_p = y_p$ (already divided by $s_p$ inside $\widetilde{\bm C}$).
  \item Compute $S_{pq}$ exactly as in~\eqref{eq:sparam}.
\end{enumerate}

\subsection{Field reconstruction}

When a full-space field is requested at $\omega$:
\begin{equation}
  \bm x(\omega) \approx \bm V\,\widetilde{\bm x}(\omega).
\end{equation}
The reduced model thus also provides VTU output at any sweep point at
post-projection cost $\mathcal{O}(nq)$. Memory permitting, only requested
frequencies need lifting; the default flag stays
\texttt{--no-write-all-fields} to keep disk usage bounded.

%---------------------------------------------------------------
\section{Cost model}

Let $n$ = full DOFs, $q$ = ROM dimension, $N_f$ = sweep points,
$N_v$ = validation points, $L$ = expansion points,
$T_{\mathrm{LU}}(n)$ = PARDISO factorization cost,
$T_{\mathrm{solve}}(n)$ = single triangular solve cost, $N_p$ = port count.

\begin{table}[h]
\centering
\caption{Asymptotic cost comparison for a sweep of $N_f$ points.}
\begin{tabular}{lll}
\toprule
Phase & Direct sweep & Krylov MOR sweep \\
\midrule
Assembly & $N_f \cdot T_{\mathrm{asm}}(n)$ & $L \cdot T_{\mathrm{asm}}(n)$ + cached subforms\\
Factorization & $N_f \cdot T_{\mathrm{LU}}(n)$ & $L \cdot T_{\mathrm{LU}}(n)$\\
Solve / Krylov build & $N_f \cdot N_p \cdot T_{\mathrm{solve}}(n)$ & $L\cdot q \cdot N_p \cdot T_{\mathrm{solve}}(n)$\\
Per-frequency online & 0 & $\mathcal{O}(N_a q^2 + q^3)$\\
Validation & 0 & $N_v \cdot T_{\mathrm{LU}}(n)$\\
\bottomrule
\end{tabular}
\end{table}

For the BP-filter benchmark used throughout this project ($n\!\approx\!2.9\!\times\!10^5$
order-1 DOFs, $N_f=401$, $L=2{-}3$, $q=30{-}60$,
$N_v\!\sim\!5$): the speedup factor is
\(
  N_f / (L + N_v) \sim 401 / 8 \approx 50,
\)
i.e.\ a $\sim$\!50$\times$ wall-clock reduction is realistic before any
parallelization.

%---------------------------------------------------------------
\section{Numerical safeguards}

\begin{itemize}
\item \textbf{Re-orthogonalization}. Modified Gram-Schmidt with one or two
      reorthogonalization passes; switch to Householder if loss of
      orthogonality exceeds $10^{-10}$.
\item \textbf{Deflation}. Drop columns whose norm after orthogonalization
      drops below $10^{-12}$; report deflation events.
\item \textbf{Conditioning}. Monitor $\|\widetilde{\bm A}^{-1}\|/\|\widetilde{\bm A}\|$ via
      LDL\textsuperscript{T} pivot magnitudes. Trigger an additional
      expansion point if conditioning exceeds $10^{12}$.
\item \textbf{Cutoff handling}. Below the lowest-order port cutoff
      $k_0 < k_{c,p}$, $\beta_p$ becomes imaginary and $s_p\to 0$. The
      reduced model must extend continuously by analytic continuation; we
      keep $\beta_p$ in the symbolic form $\sqrt{k_0^2 - k_{c,p}^2}$ and
      branch-cut on the principal sqrt. The behavior is identical to the
      existing direct solver (\texttt{powerNormalizationFactor} returns
      0).
\item \textbf{Symmetry recovery}. If round-off perturbs symmetry,
      symmetrize $\widetilde{\bm K}\!\gets\!\tfrac12(\widetilde{\bm K}+\widetilde{\bm K}^{\!\top})$
      explicitly after construction.
\end{itemize}

%---------------------------------------------------------------
\section{Mapping to existing code}

\begin{table}[h]
\centering
\begin{tabular}{p{0.36\linewidth}p{0.6\linewidth}}
\toprule
Mathematical object & Source location \\
\midrule
$\bm K - k_0^2\bm M$ & \texttt{src/fem/FEMAssembler.cpp::assemble} (volume part)\\
$\im\beta_p\bm m_p\bm m_p^{\!\top}$ & \texttt{src/fem/FEMAssembler.cpp::applyWavePorts}\\
$\bm m_p$ (output) & \texttt{src/fem/PortModeSolver.cpp::computeMode::couplingWeights}\\
$s_p(\omega)$ & \texttt{src/fem/PortModeSolver.cpp::powerNormalizationFactor}\\
$b_p$, $S_{pq}$ & \texttt{src/post/ResultExtractor.cpp::extract / portProjection}\\
PARDISO LU/LDL\textsuperscript{T} of $\bm A(\omega_0)$ & \texttt{src/linalg/MklPardisoSolver.cpp}\\
CSR pattern reuse & \texttt{src/linalg/SparseMatrix.cpp::createPattern}, \texttt{FEMAssembler::sparsityPattern\_}\\
\bottomrule
\end{tabular}
\end{table}

The MOR module will introduce a new namespace
\textcolor{gray}{\texttt{fem::mor}} living under \texttt{src/mor/} with no
upward dependency on \texttt{app}; its only inputs are the existing
\texttt{FEMAssembler}, \texttt{PortModeSolver}, and a \texttt{LinearSolver}
abstraction.

%---------------------------------------------------------------
\section{Validation strategy}

\begin{enumerate}
\item \textbf{Analytic check}. Empty rectangular waveguide with a flat
      lossless body: $S_{21}(\omega)=\mathrm{e}^{-\im\beta L}$ exactly.
      ROM with $L=2,\ q=20$ must match within $10^{-6}$ in $|S|$ over a
      2:1 band.
\item \textbf{Loaded BP-filter (project benchmark)}. The Krylov ROM with
      $L=3,\ q=40$ over $40$--$43$\,GHz vs.\ the direct 401-point sweep:
      $\rho_S \le 0.05$\,dB everywhere.
\item \textbf{Energy conservation}. For lossless cases verify
      $|S_{11}|^2+|S_{21}|^2 \le 1+10^{-3}$ at all sweep points after ROM
      evaluation (already a regression target for direct sweep).
\item \textbf{Convergence in $q$}. Plot $\max_\omega \|\widetilde{\bm
      S}-\bm S\|_F$ vs.\ $q$ at fixed $L$. Expect exponential decay until
      $q$ saturates.
\end{enumerate}

%---------------------------------------------------------------
\section{Limitations and out-of-scope items}

\begin{itemize}
\item \textbf{Strong dispersion}. The current model handles linear $\sigma$
      losses exactly and Debye/Drude models via affine expansion. Highly
      non-rational $\varepsilon(\omega)$ (e.g.\ measured tabulated data
      with sharp resonances) requires either piecewise rational fitting or
      a different reduction method (e.g.\ AAA + balanced truncation),
      which is out of scope.
\item \textbf{Strong nonlinearity}. Out of scope; non-iterative MOR is not
      meaningful here.
\item \textbf{Eigenmode tracking near band edges}. ROM reproduces the
      transfer function but does not robustly identify resonances inside
      stop-bands; a separate Krylov-Schur eigenvalue path is delegated to
      a future module.
\item \textbf{Adaptive p-refinement coupling}. ROM is built on a fixed
      mesh and basis order. h/p-adaptivity in the same sweep changes
      operators between iterations; supporting this requires re-projecting
      $\bm V$ onto the new $V_h$, which is treated in a future revision.
\end{itemize}

%---------------------------------------------------------------
\section{References (informal)}

\begin{itemize}
\item Z.\ Bai, ``Krylov subspace techniques for reduced-order modeling of large-scale dynamical systems,'' Appl.\ Numer.\ Math., 2002.
\item A.\ Odabasioglu, M.\ Celik, L.~T.\ Pileggi, ``PRIMA: passive reduced-order interconnect macromodeling algorithm,'' IEEE T-CAD, 1998.
\item J.-S.\ Zhao, J.-F.\ Lee, ``A model order reduction technique using Krylov subspace methods for the analysis of microwave circuits,'' IEEE Microw.\ Wireless Compon.\ Lett., 2002.
\item E.\ Grimme, ``Krylov projection methods for model reduction,'' PhD thesis, Univ.\ of Illinois, 1997.
\item V.\ de la Rubia, U.\ Razafison, Y.\ Maday, ``Reliable fast frequency sweep for microwave devices via the reduced basis method,'' IEEE TMTT, 2009.
\end{itemize}

\end{document}
